\chapter{Zustandsmodellierung mit State Machines}

\begin{lernziele}
Nach diesem Kapitel können Sie Betriebszustände und Zustandswechsel eines Roboters modellieren.
\end{lernziele}

\section{Zustände}
Ein Zustand beschreibt eine Situation, in der sich das System über eine gewisse Zeit befindet. Beispiele für den Roboter sind \texttt{Idle}, \texttt{Mapping}, \texttt{Navigation}, \texttt{Charging} und \texttt{Error}.

\section{Ereignisse}
Ereignisse lösen Übergänge aus. Beispiele:
\begin{itemize}
\item \texttt{GoalReceived}
\item \texttt{BatteryLow}
\item \texttt{ObstacleDetected}
\item \texttt{ChargingComplete}
\item \texttt{ErrorDetected}
\end{itemize}

\section{Beispiel}
\begin{verbatim}
Idle --GoalReceived--> Navigation
Navigation --BatteryLow--> ReturnToDock
ReturnToDock --DockReached--> Charging
Charging --ChargingComplete--> Idle
Navigation --ErrorDetected--> Error
\end{verbatim}

\section{Warum State Machines nützlich sind}
Zustandsdiagramme verhindern unklare Logik. Sie zeigen, welche Übergänge erlaubt sind. Darf der Roboter während eines Fehlers ein neues Ziel annehmen? Darf er während des Ladens navigieren? Solche Fragen werden sichtbar.

\begin{uebungbox}
Erstellen Sie eine State Machine für den Roboter mit mindestens fünf Zuständen und acht Übergängen.
\end{uebungbox}
